% Mathématiques 5e — cours V3
% Probabilités : premières expériences aléatoires

\section{Objectifs}

\begin{itemize}
\item Aborder le hasard à partir de situations simples.
\item Utiliser le vocabulaire : expérience aléatoire, issue, événement.
\item Attribuer des probabilités dans des situations simples d'équiprobabilité.
\item Positionner des événements sur une échelle de probabilité.
\item Donner une probabilité sous forme fractionnaire, décimale ou en pourcentage.
\item Relier des expressions courantes comme « une chance sur quatre » à une probabilité.
\item Répéter matériellement une expérience aléatoire simple.
\item Enregistrer des effectifs et des fréquences observés.
\item Distinguer modèle probabiliste et résultats expérimentaux.
\end{itemize}

\section{1. Une expérience aléatoire}

Une expérience aléatoire est une expérience dont on ne connaît pas à l'avance le résultat exact.

Exemples :
\begin{itemize}
\item lancer une pièce ;
\item lancer un dé ;
\item tirer une boule dans une urne ;
\item choisir au hasard une personne dans un groupe.
\end{itemize}

\begin{definition}
Une \textbf{issue} est un résultat possible d'une expérience aléatoire.
\end{definition}

\section{2. Lister les issues}

Pour un dé équilibré à six faces, les issues sont :
\[
1,\quad2,\quad3,\quad4,\quad5,\quad6.
\]

Pour une pièce :
\[
\text{pile},\quad\text{face}.
\]

Une liste complète des issues est indispensable avant tout calcul.

\coursefigure{/images/mathematiques/5eme/proba-univers.svg}{Six issues d'un dé dont les issues paires sont mises en évidence.}{Un événement peut regrouper plusieurs issues, comme « obtenir un nombre pair ».}

\section{3. Événement}

\begin{definition}
Un événement est une condition qui peut être réalisée par une ou plusieurs issues.
\end{definition}

Avec un dé :
\[
A=\text{« obtenir un nombre pair »}.
\]

Les issues favorables sont :
\[
2,\quad4,\quad6.
\]

\section{4. Événement impossible et certain}

Un événement impossible ne peut jamais se produire.

Sa probabilité vaut :
\[
0.
\]

Un événement certain se produit à chaque expérience.

Sa probabilité vaut :
\[
1.
\]

Par exemple, avec un dé classique :
\[
P(\text{obtenir }8)=0,
\]
et :
\[
P(\text{obtenir un nombre inférieur à }7)=1.
\]

\section{5. Une probabilité entre zéro et un}

Toute probabilité vérifie :
\[
\boxed{0\le P\le1}.
\]

Plus la probabilité est proche de :
\[
1,
\]
plus l'événement est probable.

Plus elle est proche de :
\[
0,
\]
moins il est probable.

\section{6. Plusieurs écritures}

Une même probabilité peut s'écrire sous différentes formes.

Par exemple :
\[
\dfrac12=0{,}5=50\%.
\]

De même :
\[
\dfrac14=0{,}25=25\%.
\]

Et :
\[
\dfrac34=0{,}75=75\%.
\]

\section{7. Langage courant}

L'expression :
\[
\text{« une chance sur quatre »}
\]
correspond à :
\[
\dfrac14.
\]

Donc :
\[
\dfrac14=25\%.
\]

De même :
\[
\text{« trois chances sur cinq »}
\]
peut être traduit par :
\[
\dfrac35.
\]

\section{8. Équiprobabilité}

Dans certaines expériences, les issues sont supposées avoir la même probabilité.

On parle d'équiprobabilité.

Un dé équilibré est modélisé par six issues équiprobables.

Une pièce équilibrée est modélisée par deux issues équiprobables.

\begin{propriete}
Dans une situation simple d'équiprobabilité :
\[
\boxed{
P(A)=
\dfrac{\text{nombre d'issues favorables}}{\text{nombre total d'issues}}
}.
\]
\end{propriete}

\section{9. Exemple : nombre pair}

Avec un dé équilibré, l'événement :
\[
A=\text{« obtenir un nombre pair »}
\]
possède trois issues favorables :
\[
2,\quad4,\quad6.
\]

Il y a six issues au total.

Donc :
\[
P(A)=\dfrac36=\dfrac12.
\]

Ainsi :
\[
\boxed{P(A)=50\%}.
\]

\section{10. Exemple : obtenir 1}

L'événement :
\[
B=\text{« obtenir }1\text{ »}
\]
possède une seule issue favorable.

Donc :
\[
P(B)=\dfrac16.
\]

Cette valeur est inférieure à :
\[
\dfrac12.
\]

L'événement est donc possible mais moins probable qu'obtenir un nombre pair.

\section{11. Urne simple}

Une urne contient :
\[
3
\]
boules rouges et :
\[
2
\]
boules bleues, indiscernables au toucher.

Il y a :
\[
5
\]
boules.

Si chaque boule a la même chance d'être tirée :
\[
P(\text{rouge})=\dfrac35,
\]
et :
\[
P(\text{bleue})=\dfrac25.
\]

On vérifie :
\[
\dfrac35+\dfrac25=1.
\]

\section{12. Hypothèse d'équiprobabilité}

L'équiprobabilité n'est pas automatique.

Si une roue comporte deux secteurs de tailles différentes, les deux couleurs ne sont pas nécessairement aussi probables.

\begin{attention}
Compter le nombre de catégories ne suffit pas à conclure à l'équiprobabilité.
\end{attention}

L'hypothèse doit être justifiée par le contexte : dé équilibré, pièce équilibrée, objets indiscernables, symétrie, etc.

\section{13. Échelle de probabilité}

On peut comparer des événements en les plaçant mentalement sur une échelle entre :
\[
0
\]
et :
\[
1.
\]

Un événement de probabilité :
\[
\dfrac16
\]
est moins probable qu'un événement de probabilité :
\[
\dfrac12.
\]

Un événement de probabilité :
\[
0{,}9
\]
est très probable mais n'est pas certain.

\section{14. Expérience répétée}

On peut répéter une même expérience plusieurs fois.

Par exemple, lancer une pièce :
\[
20
\]
fois.

On enregistre les résultats obtenus.

\begin{tabular}{c|c}
Issue & Effectif \\
Pile & $12$ \\
Face & $8$ \\
\end{tabular}

L'effectif total vaut :
\[
20.
\]

\section{15. Fréquence observée}

La fréquence observée de pile vaut :
\[
f=\dfrac{12}{20}=0{,}6.
\]

Donc :
\[
f=60\%.
\]

Pour face :
\[
f=\dfrac8{20}=0{,}4=40\%.
\]

La somme des fréquences vaut :
\[
1.
\]

\section{16. Fréquence et probabilité}

Pour une pièce équilibrée :
\[
P(\text{pile})=\dfrac12=0{,}5.
\]

Mais une série réelle peut donner :
\[
f(\text{pile})=0{,}6.
\]

Il n'y a pas de contradiction.

La probabilité décrit le modèle.

La fréquence décrit une série d'expériences réellement observée.

\coursefigure{/images/mathematiques/5eme/proba-frequences.svg}{Courbe de fréquence observée lors de répétitions comparée à une valeur théorique.}{Les fréquences fluctuent d'une série à l'autre ; elles ne reproduisent pas exactement la probabilité à chaque expérience.}

\section{17. Une petite série peut beaucoup fluctuer}

Si une pièce est lancée seulement :
\[
4
\]
fois, on peut parfaitement observer :
\[
4
\]
fois pile.

La fréquence vaut alors :
\[
1.
\]

Pourtant le modèle reste :
\[
P(\text{pile})=\dfrac12.
\]

Une courte expérience ne permet pas à elle seule de réfuter le modèle.

\section{18. Tableau d'effectifs et fréquences}

Pour une expérience répétée, un tableau peut contenir :
\begin{itemize}
\item les issues ;
\item les effectifs observés ;
\item les fréquences.
\end{itemize}

\begin{tabular}{c|c|c}
Issue & Effectif & Fréquence \\
$1$ & $8$ & $0{,}16$ \\
$2$ & $10$ & $0{,}20$ \\
$3$ & $7$ & $0{,}14$ \\
$4$ & $9$ & $0{,}18$ \\
$5$ & $8$ & $0{,}16$ \\
$6$ & $8$ & $0{,}16$ \\
\end{tabular}

L'effectif total est :
\[
50.
\]

\section{19. Comparer au modèle}

Pour un dé équilibré, chaque face a pour probabilité :
\[
\dfrac16\approx0{,}167.
\]

Les fréquences observées ne sont pas exactement égales à cette valeur, mais on peut les comparer.

Le but est de confronter :
\[
\text{modèle théorique}
\]
et :
\[
\text{observations}.
\]

\section{20. Simulation et expérience réelle}

Une expérience peut être :
\begin{itemize}
\item réalisée matériellement ;
\item simulée avec un logiciel ou un tableur.
\end{itemize}

Dans les deux cas, il faut enregistrer les résultats et conserver le nombre total d'essais.

\begin{remarque}
Une simulation informatique produit des résultats pseudo-aléatoires mais permet de répéter rapidement une expérience.
\end{remarque}

\section{21. Exemple développé}

On tire au hasard une carte parmi :
\[
10
\]
cartes numérotées de $1$ à $10$.

On cherche la probabilité d'obtenir un multiple de $3$.

Les issues favorables sont :
\[
3,\quad6,\quad9.
\]

Il y en a :
\[
3.
\]

Donc :
\[
P=\dfrac3{10}=0{,}3=30\%.
\]

Ainsi :
\[
\boxed{P=30\%}.
\]

\section{22. Méthode en équiprobabilité}

\begin{methode}
\begin{enumerate}
\item lister les issues possibles ;
\item vérifier qu'elles sont équiprobables ;
\item identifier les issues favorables ;
\item compter les favorables ;
\item diviser par le nombre total d'issues ;
\item convertir éventuellement en décimal ou pourcentage ;
\item vérifier que le résultat est compris entre $0$ et $1$.
\end{enumerate}
\end{methode}

\section{23. Méthode expérimentale}

\begin{methode}
Pour une expérience répétée :
\begin{enumerate}
\item fixer le nombre d'essais ;
\item réaliser ou simuler l'expérience ;
\item noter chaque résultat ;
\item calculer les effectifs ;
\item calculer les fréquences ;
\item comparer ces fréquences au modèle si celui-ci est connu.
\end{enumerate}
\end{methode}

\section{24. Erreurs fréquentes}

\begin{itemize}
\item Confondre issue et événement.
\item Donner une probabilité supérieure à $1$.
\item Supposer une équiprobabilité sans justification.
\item Compter deux fois une même issue.
\item Confondre fréquence observée et probabilité théorique.
\item Croire qu'une probabilité de $50\%$ impose exactement la moitié de succès dans toute série.
\item Oublier le nombre total d'essais dans le calcul d'une fréquence.
\end{itemize}

\section{À retenir}

\begin{aretenir}
Une expérience aléatoire possède plusieurs issues possibles.

Un événement regroupe une ou plusieurs issues.

Toute probabilité vérifie :
\[
\boxed{0\le P\le1}.
\]

Dans une situation d'équiprobabilité :
\[
\boxed{
P(A)=
\dfrac{\text{issues favorables}}{\text{issues possibles}}
}.
\]

Une expérience répétée fournit des fréquences :
\[
\boxed{
f=
\dfrac{\text{effectif observé}}{\text{nombre total d'essais}}
}.
\]

Probabilité et fréquence doivent être distinguées : l'une appartient au modèle, l'autre aux observations.
\end{aretenir}
